\documentclass[11pt]{article}
\usepackage[margin=1.05in]{geometry}
\usepackage{amsmath,amssymb,amsthm,mathtools}
\usepackage{microtype}
\usepackage{enumitem}
\usepackage{xcolor}
\usepackage[hidelinks]{hyperref}
\hypersetup{pdftitle={Positive Recurrence of Single-Linkage Bimolecular Weakly Reversible Stochastic Reaction Networks},pdfauthor={Alec Kriebel}}
\usepackage{booktabs}
\usepackage{longtable}

\newtheorem{theorem}{Theorem}[section]
\newtheorem{lemma}[theorem]{Lemma}
\newtheorem{proposition}[theorem]{Proposition}
\newtheorem{corollary}[theorem]{Corollary}
\theoremstyle{definition}
\newtheorem{definition}[theorem]{Definition}
\newtheorem{remark}[theorem]{Remark}

\newcommand{\N}{\mathbb N_0}
\newcommand{\Ccal}{\mathcal C}
\newcommand{\Rcal}{\mathcal R}
\newcommand{\Gamm}{\Gamma}
\newcommand{\wtG}{\widetilde\Gamma}
\newcommand{\ff}[2]{(#1)_{#2}}
\newcommand{\E}{\mathbb E}
\newcommand{\one}{\mathbf 1}

\title{Positive Recurrence of Single-Linkage Bimolecular\\
Weakly Reversible Stochastic Reaction Networks}
\author{Alec Kriebel}
\date{August 2026}

\begin{document}
\maketitle

\begin{abstract}
We prove that every finite bimolecular weakly reversible stochastic
mass-action network with one linkage class is positive recurrent on every
closed communicating class, for every positive rate vector.  No pure-species
complex assumption is imposed.  The proof augments the embedded jump chain by
the target complex of the most recent reaction and removes that target from
the population.  A residual factorial potential then has the exact increment
\(
\log(\ff{x}{t}/\ff{x}{s})
\)
when the next source is \(s\) and the carried target is \(t\).  Following the
carried target has zero reward.  A finite library of target-following episodes
reduces every possible hierarchy of source propensities to a scalar concave
recursion.  A normalized-log compactification of an arbitrary divergent
sequence then gives either an enabled source of strictly greater asymptotic
weight at a lower terminal complex, or an exact reaction-wise conservation
law contradicting divergence in the fixed class.  This yields one common
proper random-time Foster function, a finite exceptional set, and bounded
jump-count episodes.  A self-contained finite trace-chain argument converts
the certificate to finite mean return time for the original continuous-time
chain.
\end{abstract}

\tableofcontents

\section{Statement and reductions}

Let the species set be finite.  A complex is a vector \(y\in\N^d\), and a
reaction is an ordered pair \(y\to y'\).  The network is \emph{bimolecular}
when \(|y|\leq 2\) for every complex.  Its stochastic mass-action propensity
at \(x\in\N^d\) is
\[
  \lambda_{y y'}(x)=\kappa_{y y'}\ff{x}{y},
  \qquad
  \ff{x}{y}=\prod_i (x_i)_{y_i},
\]
with positive rate constants.  We use the inherited fact that every
bimolecular weakly reversible system is nonexplosive.

\begin{theorem}[Single-linkage positive recurrence]\label{thm:main}
Let a finite bimolecular stochastic mass-action reaction network have one
linkage class, and suppose its reaction graph is strongly connected.  For
every positive rate vector, every closed communicating class of its
continuous-time Markov chain is positive recurrent.
\end{theorem}

A reaction \(y\to y\) makes no state change and contributes zero to the
generator.  We discard such reactions throughout.  Duplicate reactions may
be combined by adding their rates.  If the complex graph has one vertex, all
closed classes are singletons.  If all complexes have the same molecularity,
every reaction preserves \(\sum_i x_i\), so every communicating class is
finite.  Thus the only nontrivial case has at least two complexes and an
infinite closed irreducible class \(\Gamm\).  Strong connectivity then gives
every complex a nontrivial outgoing reaction.

Write
\[
  \bar\kappa_y=\sum_{y\to u,\,u\ne y}\kappa_{yu},
  \qquad
  \bar\kappa_- = \min_y\bar\kappa_y,
  \qquad
  \bar\kappa_+ = \max_y\bar\kappa_y.
\]
All three quantities are finite and the first minimum is positive.

\section{The target-augmented jump chain}

Consider the embedded jump chain of the CTMC.  After a reaction fires, record
its target complex.  A reachable augmented state is denoted \((x,t)\), where
\(x\in\Gamm\), \(t\in\Ccal\), and \(x\geq t\).  Put
\[
  r=x-t.
\]
If the next reaction is \(s\to u\), then
\[
  (x,t)\longmapsto (x-s+u,u),
  \qquad
  r\longmapsto x-s.
\]
The transition probabilities depend on \(x\) and not on the historical label
\(t\), so the augmented process is a Markov chain.

Let \(\wtG\) be the set of augmented pairs reachable after a genuine jump.

\begin{lemma}[Augmented irreducibility]\label{lem:aug-irr}
The augmented embedded chain on \(\wtG\) is irreducible.
\end{lemma}

\begin{proof}
Take \((x,t),(x',t')\in\wtG\).  By reachability of the second pair, there is
a reaction \(s'\to t'\) and a predecessor population
\(z=x'-t'+s'\in\Gamm\) from which that reaction is enabled.  Population
irreducibility gives a finite positive-probability reaction path from \(x\)
to \(z\).  The same path can be followed from the first augmented pair,
because augmented labels do not affect propensities.  Firing
\(s'\to t'\) then reaches \((x',t')\).
\end{proof}

Define the residual factorial potential
\begin{equation}\label{eq:V}
  V(x,t)=\sum_{i=1}^d \log((x_i-t_i)!).
\end{equation}
It is nonnegative.  Since \(\log n!\to\infty\) and \(t\) belongs to a finite
set of complexes, every sublevel set of \(V\) in \(\wtG\) is finite.

\begin{lemma}[Exact target/source identity]\label{lem:residual}
At an augmented state \((x,t)\), if the next reaction has source \(s\) and
target \(u\), then
\begin{equation}\label{eq:residual-increment}
  V(x-s+u,u)-V(x,t)
  =\log\frac{\ff{x}{t}}{\ff{x}{s}}.
\end{equation}
In particular, a reaction sourced at the carried target, \(s=t\), has exactly
zero residual reward.
\end{lemma}

\begin{proof}
The new residual is \(x-s\).  Hence
\[
 \exp(\Delta V)
 =\frac{\prod_i(x_i-s_i)!}{\prod_i(x_i-t_i)!}
 =\frac{\ff{x}{t}}{\ff{x}{s}}.
\]
\end{proof}

This identity is the reason for augmenting by a target.  A product complex is
kept as a complete credit until the chain switches to a different source.

\section{A one-jump source-probability bound}

At a population state \(x\), define
\[
  \Lambda(x)=\sum_{y\in\Ccal}\bar\kappa_y\ff{x}{y},
  \qquad
  p_x(y)=\frac{\bar\kappa_y\ff{x}{y}}{\Lambda(x)}
\]
for enabled sources.  Since the carried target \(t\) is present,
\(p_x(t)>0\).  The expected one-jump residual increment is
\[
  d(x,t)=\sum_s p_x(s)\log\frac{\ff{x}{t}}{\ff{x}{s}}.
\]

\begin{lemma}[Entropy bound]\label{lem:entropy}
For every reachable augmented state,
\begin{equation}\label{eq:entropy-bound}
  d(x,t)\leq \log p_x(t)+C_0,
  \qquad
  C_0=\log|\Ccal|+
       \log\frac{\bar\kappa_+}{\bar\kappa_-}.
\end{equation}
\end{lemma}

\begin{proof}
From
\(
\ff{x}{s}=p_x(s)\Lambda(x)/\bar\kappa_s
\)
we obtain the exact identity
\[
 d(x,t)
 =\log p_x(t)-\sum_s p_x(s)\log p_x(s)
   +\sum_s p_x(s)\log\bar\kappa_s-
    \log\bar\kappa_t.
\]
The entropy is at most \(\log|\Ccal|\), and the final rate contribution is at
most \(\log(\bar\kappa_+/\bar\kappa_-)\).
\end{proof}

No comparison between independent products of rate constants is used.  The
rates enter only through the fixed additive constant \(C_0\) and positive
conditional edge probabilities below.

\section{Finite target-following episodes}

For every ordered pair \((t,c)\in\Ccal^2\), fix a simple directed path
\begin{equation}\label{eq:path}
  t=y_0\longrightarrow y_1\longrightarrow\cdots
   \longrightarrow y_L=c,
  \qquad L\leq |\Ccal|-1,
\end{equation}
which exists by strong connectivity.  Choose one particular reaction edge on
each step.

Starting from \((r+t,t)\), define the \((t,c)\)-episode as follows.  At phase
\(y_k\), \(k<L\), inspect the next jump.  Continue only if the exact
designated reaction \(y_k\to y_{k+1}\) fires; otherwise stop immediately
after that jump.  At terminal phase \(c\), take one further jump and stop.

The designated lift is
\[
  r+y_0\longrightarrow r+y_1\longrightarrow\cdots
  \longrightarrow r+c.
\]
Every designated source is literally present.  The residual stays equal to
\(r\), and closedness of \(\Gamm\) keeps every lifted state in the class.
The episode length is between one and \(|\Ccal|\) embedded jumps.

Set
\[
  q_k=\frac{\kappa_{y_k y_{k+1}}}{\bar\kappa_{y_k}}>0,
  \qquad
  p_k=p_{r+y_k}(y_k).
\]
Let \(J_k\) be the expected remaining \(V\)-reward from phase \(k\).

\begin{lemma}[Exact episode recursion]\label{lem:recursion}
The expected rewards satisfy
\begin{equation}\label{eq:recursion}
  J_L=d(r+c,c),
  \qquad
  J_k=d(r+y_k,y_k)+q_kp_kJ_{k+1}.
\end{equation}
\end{lemma}

\begin{proof}
The first jump at phase \(k\) has expected reward \(d(r+y_k,y_k)\).  Only the
exact designated edge continues.  Its probability is \(q_kp_k\), and its
immediate reward is zero by Lemma~\ref{lem:residual}.
\end{proof}

\section{Complete-credit elimination}

For \(\varepsilon\in(0,1]\), define
\[
  M_L(\varepsilon)=\log\varepsilon+C_0
\]
and, backward along the path,
\begin{equation}\label{eq:envelope}
  M_k(\varepsilon)=
  \sup_{0<p\leq1}
  \left\{\log p+C_0+q_kpM_{k+1}(\varepsilon)\right\}.
\end{equation}
Lemma~\ref{lem:entropy} and induction in
Lemma~\ref{lem:recursion} give \(J_k\leq M_k\) whenever the terminal source
probability is \(\varepsilon\).

\begin{lemma}[Scalar elimination]\label{lem:scalar}
For every fixed finite path and every fixed positive \(q_0,\ldots,q_{L-1}\),
\begin{equation}\label{eq:scalar-limit}
  M_0(\varepsilon)\longrightarrow-\infty
  \qquad\text{as }\varepsilon\downarrow0.
\end{equation}
\end{lemma}

\begin{proof}
For fixed \(M\), the function
\[
  f(p)=\log p+C_0+qpM
\]
is strictly concave.  If \(M=-A\) and \(qA>1\), its unique maximizer is
\(p=(qA)^{-1}\), and
\begin{equation}\label{eq:scalar-step}
  \sup_{0<p\leq1}f(p)
  =C_0-1-\log(qA).
\end{equation}
The terminal bound tends to minus infinity.  If
\(M_{k+1}(\varepsilon)\to-\infty\), then for all small \(\varepsilon\) the
interior formula \eqref{eq:scalar-step} applies and also tends to minus
infinity.  Backward induction proves the claim.
\end{proof}

\begin{corollary}[Complete-credit episode]\label{cor:credit}
If, along a sequence of residuals, the source probability of the terminal
complex \(c\) at \(r+c\) tends to zero, then the expected \(V\)-reward of the
entire \((t,c)\)-episode tends to minus infinity.
\end{corollary}

This finite recursion absorbs every intermediate source-rate scale.  No
averaged chain, Schur complement sign assertion, or local Lyapunov switching
is needed.

\section{Normalized-log compactification}

We now prove that every divergent sequence either gives the terminal
probability required by Corollary~\ref{cor:credit} or contradicts an exact
conservation law.

Let \((x^{(n)},t^{(n)})\) be a divergent sequence in \(\wtG\).  After taking a
subsequence, \(t^{(n)}=t\) is fixed.  Put \(r^{(n)}=x^{(n)}-t\).  By a finite
diagonal extraction, each coordinate of \(r^{(n)}\) is either constant or
tends to infinity.  Let \(I\) be the divergent coordinates and \(J=I^c\).
Define
\[
  R_n=\sum_{i\in I}\log(r_i^{(n)}+1).
\]
After a further subsequence, the limits
\begin{equation}\label{eq:w}
  w_i=\lim_{n\to\infty}
      \frac{\log(r_i^{(n)}+1)}{R_n},\quad i\in I,
  \qquad
  w_i=0,\quad i\in J,
\end{equation}
exist.  Thus \(w\geq0\) and \(\sum_iw_i=1\).  Divergent coordinates with
\(w_i=0\) encode all slower subtiers.

For \(y\in\Ccal\), put
\[
  h(y)=w\cdot y,
  \qquad
  a=\max_{y\in\Ccal}h(y),
  \qquad
  T=\{y:h(y)=a\}.
\]

\begin{lemma}[Top availability or global conservation]\label{lem:top}
Along the fixed-class divergent sequence, one of the following holds.

\begin{enumerate}[label=(\roman*)]
\item There is a real species-linear conservation law \(\ell\cdot X\) whose
      value tends to infinity along the sequence, contradicting membership
      in one communicating class.
\item There are complexes \(s,c\in\Ccal\) with \(h(s)>h(c)\) such that, for
      every sufficiently large \(n\), \(s\) is enabled at \(r^{(n)}+c\).
\end{enumerate}
\end{lemma}

\begin{proof}
If every complex is in \(T\), then
\(w\cdot(y'-y)=0\) for every reaction.  Hence \(w\cdot X\) is conserved.
Some positive-weight coordinate diverges, giving (i).

Assume \(T\ne\Ccal\).  If a top complex \(s\) contains two particles from
\(I\), bimolecularity says it contains no bounded-coordinate particle.  It
is enabled at \(r^{(n)}+c\) for every lower complex \(c\notin T\), giving
(ii).

It remains to suppose no top complex contains two \(I\)-particles.  Since the
maximum is positive, every top complex contains exactly one \(I\)-particle.
Let \(K\subseteq I\) be the species appearing in top complexes.  Each
\(i\in K\) has \(w_i=a\).  No complex contains two \(K\)-particles, and
\begin{equation}\label{eq:top-qk}
  y\in T\quad\Longleftrightarrow\quad
  q_K(y):=\sum_{i\in K}y_i=1.
\end{equation}
Indeed, a complex containing one \(K\)-particle already has weight at least
\(a\), so it must have exactly weight \(a\); two would have weight \(2a>a\).
Conversely every top complex contains a species in \(K\).

If every complex has \(q_K=1\), then \(M_K=\sum_{i\in K}X_i\) is conserved
and diverges, giving (i).  Otherwise there is a lower complex with
\(q_K=0\).

If a unary top complex \(K_i\) exists, it is enabled over any lower terminal
complex, giving (ii).  Otherwise each top complex has the form \(K_i+D\).
The service species \(D\) lies in \(J\): if \(D\in I\), the top complex would
contain two \(I\)-particles, contrary to the current case.

If a lower complex \(c\) contains a service species \(D\), the corresponding
source \(K_i+D\) is enabled at \(r^{(n)}+c\), and (ii) follows.  If no lower
complex contains any service species, then
\begin{equation}\label{eq:token-law}
  M_K-\sum_{D\text{ service}}X_D
\end{equation}
is zero on every complex: a top complex contains one \(K\)-particle and one
service particle, while a lower complex contains neither.  It is therefore
an exact reaction-wise conservation law.  Its positive part diverges and its
negative-coordinate species all belong to \(J\), hence remain fixed along
the subsequence.  This gives (i).
\end{proof}

The conservation laws in the proof are global on the entire linkage class;
no local phase coboundary is promoted to a species invariant.

\begin{lemma}[Vanishing terminal probability]\label{lem:terminal}
Alternative (ii) of Lemma~\ref{lem:top} implies
\[
  p_{r^{(n)}+c}(c)\longrightarrow0.
\]
\end{lemma}

\begin{proof}
For an enabled fixed bimolecular complex \(y\), coordinatewise comparison
with \eqref{eq:w} gives
\begin{equation}\label{eq:falling-asymptotic}
  \frac{1}{R_n}\log\ff{r^{(n)}+c}{y}
  \longrightarrow h(y).
\end{equation}
Positive-weight divergent coordinates give their normalized logarithms;
zero-weight divergent coordinates contribute \(o(R_n)\); fixed coordinates
contribute \(O(1)\).  Both \(s\) and \(c\) are enabled.  Since
\(h(s)>h(c)\),
\[
  \frac{\ff{r^{(n)}+c}{s}}
       {\ff{r^{(n)}+c}{c}}
  \longrightarrow\infty.
\]
Therefore
\[
 p_{r^{(n)}+c}(c)
 \leq
 \frac{\bar\kappa_c\ff{r^{(n)}+c}{c}}
      {\bar\kappa_s\ff{r^{(n)}+c}{s}}
 \longrightarrow0.
\]
\end{proof}

\section{Finite-family uniformization}

For each \(c\in\Ccal\), let \(D_c(x,t)\) be the expected \(V\)-increment of
the fixed \((t,c)\)-episode.  Define
\begin{equation}\label{eq:K}
  K=\left\{(x,t)\in\wtG:
       \min_{c\in\Ccal}D_c(x,t)>-1\right\}.
\end{equation}

\begin{proposition}[Finite-flag uniformization]\label{prop:uniform}
The set \(K\) is finite.
\end{proposition}

\begin{proof}
If \(K\) were infinite, properness of \(V\) would give a divergent sequence
in \(K\).  Apply the compactification above.  A conservation alternative in
Lemma~\ref{lem:top} contradicts the fixed communicating class.  Otherwise,
Lemmas~\ref{lem:terminal} and \ref{lem:scalar} give a fixed terminal complex
\(c\) such that
\[
  D_c(x^{(n)},t)\longrightarrow-\infty.
\]
This contradicts \((x^{(n)},t)\in K\), where every terminal episode has drift
greater than \(-1\).
\end{proof}

Outside \(K\), choose before the episode starts the lexicographically first
terminal complex minimizing \(D_c(x,t)\).  This is a deterministic measurable
selector on the countable state space.  It is not chosen after observing an
outcome.  The selected stopping time \(\tau\) satisfies
\begin{equation}\label{eq:foster}
  1\leq\tau\leq L:=|\Ccal|,
  \qquad
  \E_{(x,t)}[V(Z_\tau)-V(x,t)]\leq-1,
  \quad (x,t)\notin K.
\end{equation}

There is no defect-promotion gap.  The episode has at most \(L\) jumps, so
every population coordinate changes by at most \(2L\).  The rank is simply
the number of designated path edges remaining: a designated jump decreases
it, and every other jump stops the episode.

\section{The random-time Foster conclusion}

We prove the needed countable-chain statement directly.

Let \(Y_0=z\) and let \(Y_{n+1}\) be the endpoint of the episode selected at
\(Y_n\).  Inside \(K\), define an episode to be one ordinary jump.  Put
\[
  \sigma_K=\inf\{n\geq0:Y_n\in K\}.
\]
For every integer \(N\), summing the conditional drift inequalities up to the
bounded index \(N\wedge\sigma_K\) gives
\begin{equation}\label{eq:episode-hit}
  \E_z V(Y_{N\wedge\sigma_K})
  +\E_z(N\wedge\sigma_K)
  \leq V(z).
\end{equation}
Since \(V\geq0\), monotone convergence yields
\begin{equation}\label{eq:mean-episodes}
  \E_z\sigma_K\leq V(z).
\end{equation}
The augmented embedded chain reaches \(K\) in at most \(L\sigma_K\) jumps,
so its expected hitting time of \(K\) is finite from every state.

We now close the finite-set-to-single-state step explicitly.

\begin{proposition}[Finite trace-chain closure]\label{prop:trace}
The augmented embedded chain has a state with finite expected return time.
\end{proposition}

\begin{proof}
For \(k\in K\), take one ordinary jump.  There are finitely many possible
successors, each with finite expected hitting time of \(K\).  Hence the next
return time to \(K\) has finite expectation from every \(k\in K\).

Record successive visits to \(K\).  The resulting trace chain is finite and
irreducible: any augmented path between two points of \(K\), with its
intermediate \(K\)-visits retained, gives a positive trace path.  Fix
\(k_*\in K\).  Finiteness and irreducibility give integers \(m\) and
\(\epsilon>0\) such that, from every trace state, \(k_*\) is reached within
\(m\) trace transitions with probability at least \(\epsilon\).  Geometric
blocking gives finite expected trace return time to \(k_*\).

Finally set
\[
  B=\max_{k\in K}\E_k T_K^+<\infty.
\]
Strong Markov conditioning bounds the expected number of original jumps in a
random collection of trace excursions by \(B\) times the expected number of
excursions.  The augmented embedded return time to \(k_*\) is therefore
finite.
\end{proof}

Projecting away the target label shows that the original population embedded
chain has finite expected return time to the population coordinate of
\(k_*\).

\section{Conversion to physical time and proof of the theorem}

Let
\[
  \kappa_*=
  \min\{\kappa_{yu}:y\ne u\}>0.
\]
At every state of an infinite irreducible class, some nontrivial reaction is
enabled.  Its falling-factorial factor is a positive integer and hence at
least one.  Therefore the total CTMC rate is at least \(\kappa_*\), and every
conditional mean holding time is at most \(1/\kappa_*\).

If \(N\) is the finite-mean embedded return time found above and \(H_j\) are
the physical holding times, Tonelli's theorem gives
\[
  \E\sum_{j=0}^{N-1}H_j
  =\E\sum_{j\geq0}\one_{\{j<N\}}
      \frac{1}{\Lambda(X_j)}
  \leq \frac{\E N}{\kappa_*}<\infty.
\]
Nonexplosion excludes accumulation of jumps.  Thus one state of \(\Gamm\)
has finite expected CTMC return time.  Irreducibility gives the same
classification for every state, proving Theorem~\ref{thm:main}.

For completeness, every selected episode has at most \(L\) jumps and each
live phase has total rate at least \(\kappa_*\).  Its physical duration is
stochastically dominated by a Gamma random variable with shape \(L\) and
rate \(\kappa_*\), so all duration moments are uniformly finite.  The
coordinate overshoot is deterministically at most \(2L\).

\section{Scope and verification}

The proof automatically respects coordinate faces, siphons, parity classes,
and other lattice restrictions because every constructed path consists of
actual enabled reactions lifted inside the given closed class.  No
irreducibility of the whole orthant is assumed.

The one-linkage hypothesis is used in exactly one global place: for every
current target and every terminal complex, the finite directed path
\eqref{eq:path} exists.  The argument therefore makes no claim for arbitrary
multiple-linkage networks.

The accompanying deterministic verifier checks:

\begin{itemize}[leftmargin=2em]
\item the exact falling-factorial residual identity on small states;
\item the entropy/source-probability rewrite;
\item equality of the episode recursion and independent branch enumeration;
\item the sharp scalar-envelope calculus;
\item the complete top availability/conservation case split;
\item an exhaustive atlas of 57,288 three-species complex-set/weight cases;
\item 20,000 independent random four-species cases;
\item canonical trigger-and-drain, immigration-death, and multi-type
      calibrations.
\end{itemize}

None of the finite enumeration is used as a universal proof.  The universal
steps are the exact identities and Lemma~\ref{lem:top}.

\section*{Research-process note}
The proof and verification package were developed in an autonomous
AI-assisted research workflow.  The mathematical claim is presented only
through arguments and deterministic checks that can be audited independently.

\appendix
\section{Audit against common invalid arguments}

The proof does not infer recurrence from nonexplosion, does not normalize an
arbitrary stationary measure, does not use deterministic persistence, does
not assert one-step negativity of total count or ordinary entropy, and does
not invoke complex balance.  Reachability alone is not treated as drift: all
competing jumps enter the exact finite recursion.  The embedded-to-continuous
conversion is explicit, and boundary states are handled by the carried-target
lift.

The Phase-IV bounded-defect computations are preserved but are not used as a
black box.  The final target-following certificate explicitly tracks target
retention, population feasibility, path containment, positive continuation
probabilities, deterministic overshoot, and physical duration.

\end{document}
